\documentclass[12pt]{article}
\usepackage[T1]{fontenc}
\usepackage[utf8]{inputenc}
\usepackage{lmodern}
\usepackage{textcomp}
\usepackage{enumitem}
\usepackage{geometry}
\geometry{margin=1in}
\begin{document}
\begin{center}
Answer Key - Astronomy - Graduate Level Assessment 4 \
\small (Answers are ordered exactly as in the question paper.)
\end{center}

\section*{Multiple Choice}
\begin{enumerate}[label=\arabic*.]
\item \textbf{Answer:} A
\\ \textit{Options:} A) asteroids would orbit with a period half that of Jupiter; B) asteroids would orbit with a period twice that of Jupiter; C) asteroids would orbit with a period twice that of Mars; D) A and B
\\ \textit{Note:} Correct option: A - asteroids would orbit with a period half that of Jupiter
\item \textbf{Answer:} D
\\ \textit{Options:} A) fast rotation only; B) a rocky mantle only; C) a molten metallic core only; D) both a molten metallic core and reasonably fast rotation
\\ \textit{Note:} Correct option: D - both a molten metallic core and reasonably fast rotation
\item \textbf{Answer:} C
\\ \textit{Options:} A) a New Horizons flyby in the 1990 s; B) Hubble Space Telescope images that resolved Pluto's disk; C) brightness measurements made during mutual eclipses of Pluto and Charon; D) radar observations made by the Arecibo telescope
\\ \textit{Note:} Correct option: C - brightness measurements made during mutual eclipses of Pluto and Charon
\item \textbf{Answer:} C
\\ \textit{Options:} A) a composition dominated by hydrogen and helium; B) the presence of belts zones and storms; C) an equatorial wind speed of more than 900 miles per hour; D) significant "shear" between bands of circulation at different latitudes
\\ \textit{Note:} Correct option: C - an equatorial wind speed of more than 900 miles per hour
\item \textbf{Answer:} C
\\ \textit{Options:} A) the entire planets are made mostly of metal.; B) metals condensed first in the solar nebula and the rocks then accreted around them.; C) metals sank to the center during a time when the interiors were molten throughout.; D) radioactivity created metals in the core from the decay of uranium.
\\ \textit{Note:} Correct option: C - metals sank to the center during a time when the interiors were molten throughout.
\end{enumerate}

\section*{True / False}
\begin{enumerate}[label=\arabic*.]
\setcounter{enumi}{5}
\item \textbf{Answer:} False
\\ \textit{Options:} A) True; B) False
\\ \textit{Note:} LLM-generated false statement. Correct answer: '11'.
\item \textbf{Answer:} False
\\ \textit{Options:} A) True; B) False
\\ \textit{Note:} LLM-generated false statement. Correct answer: 'the rigid rocky material of the crust and uppermost portion of the mantle.'.
\item \textbf{Answer:} False
\\ \textit{Options:} A) True; B) False
\\ \textit{Note:} LLM-generated false statement. Correct answer: 'You see the star's blackbody spectrum with absorption lines due to hydrogen.'.
\item \textbf{Answer:} True
\\ \textit{Options:} A) True; B) False
\\ \textit{Note:} LLM-generated true statement based on MCQ.
\item \textbf{Answer:} True
\\ \textit{Options:} A) True; B) False
\\ \textit{Note:} LLM-generated true statement based on MCQ.
\end{enumerate}

\section*{Long Form Response}
\begin{enumerate}[label=\arabic*.]
\setcounter{enumi}{10}
\item \textbf{Answer:} Jupiter's greater mass compresses it more thus increasing its density.. Explain why this is the correct answer and provide supporting evidence. Reference key facts or concepts that support this choice.
\\ \textit{Note:} Long-form derived from MCQ; award credit for stating the correct idea and giving supporting reasoning.
\item \textbf{Answer:} gamma rays X rays ultraviolet visible light infrared radio. Explain why this is the correct answer and provide supporting evidence. Reference key facts or concepts that support this choice.
\\ \textit{Note:} Long-form derived from MCQ; award credit for stating the correct idea and giving supporting reasoning.
\end{enumerate}

\vfill
\noindent\textit{End of Answer Key}
\end{document}